\section{Introduction}

\IEEEPARstart{L}{oRa} and LoRaWAN have become important low-power wide-area networking technologies for smart-city sensing, industrial monitoring, environmental data collection, and other large-scale IoT deployments.
In these scenarios, end devices are often inexpensive, battery-powered, and deployed for long periods in uncontrolled environments.
Although upper-layer cryptographic mechanisms are essential, they may not fully eliminate risks caused by cloned devices, compromised credentials, rogue transmitters, or post-deployment hardware replacement.
This motivates lightweight edge-side authentication mechanisms that can complement protocol-layer security without requiring additional hardware on constrained IoT nodes~\cite{lorawan_spec,hamdaoui2022deeplearning}.

Radio frequency fingerprint identification (RFFI) addresses this need by exploiting device-specific transmitter impairments introduced by analog front-end components, including power-amplifier nonlinearity, I/Q imbalance, oscillator imperfections, and filtering effects.
These impairments can leave subtle but repeatable signatures in received baseband signals, enabling a gateway to identify or verify a transmitter from its RF emissions~\cite{brik2008wireless_device_id,jian2020deeplearning_rf_fingerprinting}.
For LoRa devices, however, practical RFFI is complicated by chirp spread spectrum modulation and by deployment variability.
A model trained under one condition may degrade under cross-day channel drift, distance changes, location changes, configuration shifts, or receiver hardware mismatch~\cite{elmaghbub2021lora_wild,hamdaoui2022deeplearning}.
Therefore, a useful LoRa RFFI model should not only perform well under a fixed collection setting, but should also remain stable under realistic deployment shifts.

Existing LoRa RFFI methods have explored raw-IQ convolutional baselines, spectrogram-based classifiers, and the use of out-of-band (OOB) distortion as additional evidence of hardware impairments~\cite{shen2021lora_spectrogram_cnn,elmaghbub2020widescan,hamdaoui2022deeplearning}.
These studies show that OOB spectrum can contain useful transmitter information; therefore, our work does not claim the first use of OOB evidence.
Instead, we focus on a different question: how should OOB evidence be fused into an RF representation so that it improves robustness rather than introducing unstable shortcuts?
Our experiments show that naive OOB concatenation and gated fusion can collapse under small-sample cross-day evaluation, whereas cross-attentive OOB fusion provides a more stable mechanism for injecting spectral hardware evidence into RF token representations.

In this paper, we propose \model, an OOB-guided cross-attentive RF-HSTU hybrid model for LoRa device authentication.
The model uses a lightweight CNN stem to preserve local RF impairment modeling, converts received LoRa captures into patch-level RF tokens, and applies an RF-HSTU sequence encoder to model token interactions.
An auxiliary OOB branch extracts spectral distortion evidence, which is then injected through cross-attention instead of static feature concatenation.
Chirp-aware positional embeddings are included as a LoRa-structure prior, while file-level mean-logits voting is used to match authentication-style evaluation where each device contributes one test capture file.

The main contributions of this paper are summarized as follows:
\begin{itemize}
  \item We propose an OOB-guided cross-attentive RF-HSTU hybrid architecture for LoRa RFFI, combining local RF distortion extraction, RF token sequence modeling, and OOB spectral evidence.
  \item We demonstrate through controlled fusion/chirp ablations that cross-attentive OOB fusion is the key mechanism for stable cross-day generalization, while naive concatenation and gated OOB fusion can be unstable under the same protocol.
  \item We conduct a strict source-domain validation evaluation on the public OSU LoRa dataset.
        Under the Day1--3/Day4/Day5 cross-day protocol, the proposed model improves file-level accuracy from 54.2$\pm$14.2\% for CNN-IQ to 75.0$\pm$5.3\%, with a paired bootstrap confidence interval entirely above zero.
  \item We further evaluate deployment shifts and receiver shifts.
        The hybrid model shows promising gains under distance and indoor-location shifts, while strict source-only cross-receiver transfer remains close to chance, revealing receiver calibration mismatch as an open challenge rather than a solved problem.
\end{itemize}

The rest of this paper is organized as follows.
Section~\ref{sec:related} reviews LoRa RFFI, OOB fingerprinting, RF sequence modeling, and domain-shift-related work.
Section~\ref{sec:system} formulates the LoRa RFFI authentication problem.
Section~\ref{sec:method} presents the proposed OOB-guided cross-attentive RF-HSTU model.
Section~\ref{sec:experiments} describes the dataset, protocols, baselines, and implementation details.
Section~\ref{sec:results} reports cross-day, ablation, deployment-shift, edge-deployment, and cross-receiver stress-test results.
Section~\ref{sec:discussion} discusses deployment implications and limitations, and Section~\ref{sec:conclusion} concludes the paper.
